\begin{lemma}[Chain rule for the weak Euclidean gradient]
Fix $n \in \mathbb{N}_0$ and let $\gamma_n$ be the Gaussian measure of
\noderef{GaussMeasure}.  Let $q, q' : \mathbb{R} \to \mathbb{R}$ and let
$L_q \in \mathbb{R}$ satisfy $L_q \geq 0$, and assume
\begin{enumerate}
  \item for every $t \in \mathbb{R}$ the function $q$ is differentiable at $t$
    with derivative $q'(t)$;
  \item $q'$ is continuous (so $q \in C^1(\mathbb{R})$);
  \item $|q'(t)| \leq 1$ for every $t \in \mathbb{R}$;
  \item $q$ is $1$-Lipschitz: $|q(s) - q(t)| \leq |s - t|$ for all $s,t$;
  \item $q'$ is $L_q$-Lipschitz: $|q'(s) - q'(t)| \leq L_q\,|s - t|$ for all $s,t$;
  \item $q(0) = 0$.
\end{enumerate}
Let further $U : \mathbb{R}^n \to \mathbb{R}$ and
$G = (G_1,\ldots,G_n) : \mathbb{R}^n \to \mathbb{R}^n$ satisfy
\begin{enumerate}
  \setcounter{enumi}{6}
  \item $U \in L^2(\gamma_n)$;
  \item each coordinate function $x \mapsto G_j(x)$ lies in $L^2(\gamma_n)$,
    for every $j \in \{1,\ldots,n\}$;
  \item $G$ is a weak Euclidean gradient of $U$ in the sense of
    \noderef{HasGaussWeakGrad}.
\end{enumerate}
Then the composition $x \mapsto q(U(x))$ has weak Euclidean gradient
$x \mapsto q'(U(x))\,G(x)$, again in the sense of \noderef{HasGaussWeakGrad};
here $q'(U(x))\,G(x)$ is the scalar multiple of the vector $G(x) \in \mathbb{R}^n$
by the real number $q'(U(x))$, so its $j$-th coordinate is
$q'(U(x))\,G_j(x)$.  Explicitly: for every $j \in \{1,\ldots,n\}$ and every test
function $\varphi \in C_c^\infty(\mathbb{R}^n)$,
\[
  \int_{\mathbb{R}^n} q(U(x))\, \partial_j \varphi(x)\, dx
    = -\int_{\mathbb{R}^n} q'(U(x))\,G_j(x)\, \varphi(x)\, dx ,
\]
both integrals being taken with respect to Lebesgue measure.
\end{lemma}
\begin{proof}
Fix $j \in \{1,\ldots,n\}$ and $\varphi \in C_c^\infty(\mathbb{R}^n)$; we must
prove the displayed identity.

Step 0 (consequences of the hypotheses on $q$).  Differentiability at every
point gives continuity of $q$.  Taking $t = 0$ in the $1$-Lipschitz hypothesis
and using $q(0) = 0$ gives
\[
  |q(t)| = |q(t) - q(0)| \leq |t| \qquad (t \in \mathbb{R}), \tag{Q}
\]
so $q$ grows at most linearly; this is what makes $q \circ U$ as integrable as
$U$ itself.

Step 1 (passing to Lebesgue measure on a compact set).  By
\noderef{GaussL2LocallyIntegrable}, a function in $L^2(\gamma_n)$ is locally
Lebesgue-integrable and lies in $L^2(K,\lambda)$ for every compact
$K \subseteq \mathbb{R}^n$, $\lambda$ denoting Lebesgue measure; the reason is
that the Gaussian density of \noderef{GaussMeasure} is continuous and strictly
positive, hence bounded below on $K$.  Applying this to $U$ and to $G_j$ gives
$U, G_j \in L^2(K,\lambda)$ for every compact $K$.  Put
\[
  K = \operatorname{supp}\varphi \cup \operatorname{supp}(\partial_j\varphi),
\]
a compact set, and note that $\varphi$ and $\partial_j\varphi$ both vanish off
$K$ and are bounded, say $|\varphi| \leq M_\varphi$.  Consequently every integral
appearing below may be restricted to $K$, and each is finite: by (Q) and
Cauchy--Schwarz on $(K,\lambda)$, $q(U)\,\partial_j\varphi$ is integrable
because $|q(U)| \leq |U|$ with $U, \partial_j\varphi \in L^2(K,\lambda)$, and
$q'(U)\,G_j\,\varphi$ is integrable because $|q'(U)| \leq 1$, $\varphi$ is
bounded and $G_j \in L^1(K,\lambda)$.

Step 2 (mollification).  Fix a mollifier family $(\rho_\varepsilon)_{\varepsilon>0}$
as in \noderef{MollifierL2Local}: each $\rho_\varepsilon$ is smooth, nonnegative,
supported in the closed ball of radius $\varepsilon$, with
$\int \rho_\varepsilon \, d\lambda = 1$.  Set
\[
  U_\varepsilon = \rho_\varepsilon * U, \qquad
  G_{j,\varepsilon} = \rho_\varepsilon * G_j .
\]
Since $U$ is locally integrable and $\rho_\varepsilon$ is smooth with compact
support, $U_\varepsilon$ is smooth and $G_{j,\varepsilon}$ is continuous.

Step 3 (mollification commutes with the weak gradient).  For every $x$,
\[
  \partial_j U_\varepsilon(x) = G_{j,\varepsilon}(x) .
\]
Indeed, differentiating under the convolution integral is legitimate here
($\rho_\varepsilon$ is smooth and compactly supported, $U$ locally integrable),
and it turns $\partial_j(\rho_\varepsilon * U)(x)$ into the pairing of $U$
against $\partial_j$ of the translated and reflected bump
$t \mapsto \rho_\varepsilon(x-t)$, which is itself a $C_c^\infty$ test function.
Hypothesis (9) applied to that test function replaces the pairing with $U$ by the
pairing with $G_j$, which is exactly $(\rho_\varepsilon * G_j)(x)$.

Step 4 (the identity for the mollified function).  $U_\varepsilon$ is smooth and
$q \in C^1$, so the classical chain rule gives
$\partial_j (q \circ U_\varepsilon) = q'(U_\varepsilon)\,\partial_j U_\varepsilon$
pointwise, and all the functions involved are continuous.  Integration by parts
against the compactly supported $\varphi$, with no boundary term, then yields
\[
  \int_{\mathbb{R}^n} q(U_\varepsilon(x))\, \partial_j \varphi(x)\, dx
    = -\int_{\mathbb{R}^n} q'(U_\varepsilon(x))\, \partial_j U_\varepsilon(x)\,
        \varphi(x)\, dx
    = -\int_{\mathbb{R}^n} q'(U_\varepsilon(x))\, G_{j,\varepsilon}(x)\,
        \varphi(x)\, dx ,
\]
the last step by Step 3.  This is the asserted identity with
$(U, G_j)$ replaced by $(U_\varepsilon, G_{j,\varepsilon})$.

Step 5 (the two $L^2(K)$ convergences).  By \noderef{MollifierL2Local} applied
to $U$ and to $G_j$, both of which lie in $L^2(K',\lambda)$ for every compact
$K'$ by Step 1,
\[
  \|U_\varepsilon - U\|_{L^2(K,\lambda)} \to 0,
  \qquad
  \|G_{j,\varepsilon} - G_j\|_{L^2(K,\lambda)} \to 0
  \qquad (\varepsilon \downarrow 0). \tag{C}
\]
In particular $\|G_{j,\varepsilon}\|_{L^2(K,\lambda)}$ is bounded for small
$\varepsilon$, being convergent.

Step 6 (passing to the limit on the left).  Using the $1$-Lipschitz hypothesis
pointwise and then Cauchy--Schwarz on $(K,\lambda)$,
\[
  \Bigl| \int_K \bigl(q(U_\varepsilon) - q(U)\bigr) \partial_j\varphi \, d\lambda \Bigr|
    \leq \int_K |U_\varepsilon - U|\, |\partial_j\varphi| \, d\lambda
    \leq \|U_\varepsilon - U\|_{L^2(K,\lambda)}\,
         \|\partial_j\varphi\|_{L^2(K,\lambda)} ,
\]
which tends to $0$ by (C), the second factor being finite since
$\partial_j\varphi$ is bounded with support in $K$ and $\lambda(K) < \infty$.
Hence the left-hand side of Step 4 converges to
$\int q(U)\,\partial_j\varphi \, d\lambda$.

Step 7 (passing to the limit on the right).  Split the difference of the
integrands as
\[
  q'(U_\varepsilon) G_{j,\varepsilon} \varphi - q'(U) G_j \varphi
    = \bigl(q'(U_\varepsilon) - q'(U)\bigr) G_{j,\varepsilon} \varphi
      + q'(U) \bigl(G_{j,\varepsilon} - G_j\bigr) \varphi .
\]
For the first term, the $L_q$-Lipschitz hypothesis on $q'$ gives
$|q'(U_\varepsilon) - q'(U)| \leq L_q |U_\varepsilon - U|$ pointwise, so with
$|\varphi| \leq M_\varphi$ and Cauchy--Schwarz on $(K,\lambda)$,
\[
  \Bigl| \int_K \bigl(q'(U_\varepsilon) - q'(U)\bigr) G_{j,\varepsilon} \varphi
    \, d\lambda \Bigr|
    \leq L_q M_\varphi \, \|U_\varepsilon - U\|_{L^2(K,\lambda)}\,
          \|G_{j,\varepsilon}\|_{L^2(K,\lambda)} ,
\]
which tends to $0$ by (C) together with the boundedness of
$\|G_{j,\varepsilon}\|_{L^2(K,\lambda)}$ noted in Step 5.  For the second term,
$|q'(U)| \leq 1$ by hypothesis (3), so again by Cauchy--Schwarz
\[
  \Bigl| \int_K q'(U) \bigl(G_{j,\varepsilon} - G_j\bigr) \varphi \, d\lambda
  \Bigr|
    \leq M_\varphi \, \|G_{j,\varepsilon} - G_j\|_{L^2(K,\lambda)}\,
          \lambda(K)^{1/2} ,
\]
which tends to $0$ by (C).  Hence the right-hand side of Step 4 converges to
$-\int q'(U)\,G_j\,\varphi \, d\lambda$.

Step 8 (conclusion).  For each $\varepsilon > 0$ the two sides of Step 4 are
equal, and by Steps 6 and 7 they converge respectively to
$\int q(U)\,\partial_j\varphi \, d\lambda$ and
$-\int q'(U)\,G_j\,\varphi\, d\lambda$ as $\varepsilon \downarrow 0$.  A
convergent net of reals has a unique limit, so the two limits agree:
\[
  \int_{\mathbb{R}^n} q(U(x))\, \partial_j \varphi(x)\, dx
    = -\int_{\mathbb{R}^n} q'(U(x))\,G_j(x)\, \varphi(x)\, dx .
\]
Since $j$ and $\varphi$ were arbitrary, and the $j$-th coordinate of
$q'(U(x))\,G(x)$ is $q'(U(x))\,G_j(x)$, the function $x \mapsto q(U(x))$ has
weak Euclidean gradient $x \mapsto q'(U(x))\,G(x)$, as asserted.
\end{proof}
